\begin{abstract}
Predicting the Remaining Useful Life (RUL) of aircraft engines enables
condition-based maintenance and reduces unplanned downtime.  This work
addresses the problem on the NASA C-MAPSS FD001 benchmark, which provides
multivariate sensor time series from 100 training and 100 test turbofan
engines run to failure.  In Phase~1 we established baselines with
engineered rolling-window features and four regression models (OLS,
Elastic Net, Random Forest, Gradient Boosting), achieving a best
validation MAE of 25.03 cycles (Random Forest).  Phase~2 replaces the
hand-crafted pipeline with end-to-end deep learning: LSTM, GRU, 1D-CNN
and Transformer architectures trained on 30-cycle sliding windows, plus a
hybrid stacking model that combines deep-learning embeddings with
traditional features.  Systematic ablation studies over window length,
input features, RUL target capping, model capacity, regularisation and
architecture quantify each design choice.  The best single deep-learning
model (LSTM) reduces the validation MAE to 19.41 cycles, a 22\%
improvement over the Phase~1 baseline, while the stacking hybrid achieves
21.48 cycles with more stable predictions across RUL bands.  All results
are reproducible with seed~42 and all code is publicly available.
\end{abstract}
